% !TeX encoding = UTF-8
% !TeX spellcheck = en_EN-EnglishUnitedKingdom
%% =======================================================
%%  P3CTeX — Example document (LaTeX2e API showcase)
%%  Build from tex/tests/ so that img/ finds tex/tests/img/:
%%    set TEXINPUTS=.;../latex;../code
%%    pdflatex P3CTeX-example.tex
%%  Or from repo root with --include-directory for tex/latex and tex/code,
%%  and run from tex/tests/ (or adjust image paths).
%% =======================================================

\documentclass[fake,cover-config=offset,default]{P3CTeX}

\PECTeXconfig{
	cover-config={
		text-minititle = {P3CTeX API Example},
		text-author    = {Author Name}
	},
	page-config={
		doc-name  = {P3CTeX Example},
		subj-name = {Software Engineering},
		refcode   = {EXAMPLE-001}
	},
	pdf-config={
		pdf-title  = {P3CTeX LaTeX2e API Example},
		pdf-author = {P3CTeX User}
	},
	cover,
	toc,
	default
}


\pxSRCsetup{placeholder-color=blue}
\begin{document}

% ========== Cover and TOC (pxGDX) ==========
%\PECTeXprintCover
%\PECTeXprintTOC*

% ========== 1. pxGDX — Sections and rules ==========
\sectionuoc{First UOC section}
Sample text with UOC section style.

\subsectionuoc{UOC subsection}
\subsubsectionuoc{UOC subsubsection}
More text with UOC subsection style.

Thin horizontal rule:\\
\pxHRule

Thick horizontal rule (centered):\\
\pxPageRule*

% ========== 2. pxPRP — Property maps ==========
\section{pxPRP: property maps and lists}

\subsection{Objects, keys, and access}
\pxNewObject{DEMO}
\pxSetValue{DEMO}{k1}{v1,v2}
\pxSetValue{DEMO}{k2}{vA,vB,vC}

\verb*|\pxGET DEMO.k1[2]| $\to$ \pxGET DEMO.k1[2].\\
\verb*|\pxPRINTF DEMO.k2[2]| $\to$ \pxPRINTF DEMO.k2[2].

\subsection{Overwrite and append}
\pxNewObject{OVER}
\pxSetValue{OVER}{state}{old}
Before overwrite: OVER.state[1] = \pxGET OVER.state[1].
\pxSetValue{OVER}{state}{newA,newB}
After: OVER.state[1] = \pxGET OVER.state[1],\quad
OVER.state[2] = \pxGET OVER.state[2].

\subsection{Lists: \texttt{\textbackslash pxPutValue}, \texttt{\textbackslash pxListSet}, \texttt{\textbackslash pxListAll}}
\pxNewObject{LST}
\pxPutValue{LST}{items}{a}
\pxPutValue{LST}{items}{b,c}
LST.items = (\pxGET LST.items[1], \pxGET LST.items[2], \pxGET LST.items[3]).

\pxNewObject{OBJL}
\pxListSet{OBJL}{nums}{1,2,3}
All values (expected 1--3): \pxListAll{OBJL}{nums}[ -- ].

\subsection{Clear key/object and predicates}
\pxNewObject{TMP}
\pxSetValue{TMP}{k}{X,Y}
Before clear: TMP.k[1] = \pxGET TMP.k[1].
\pxClearKey{TMP}{k}
\pxSetValue{TMP}{k}{Z}
After: TMP.k[1] = \pxGET TMP.k[1].
\pxClearObject{TMP}
\pxNewObject{TMP}
\pxSetValue{TMP}{k}{P,Q}
Now: TMP.k[1] = \pxGET TMP.k[1].

\pxIfObjectExistsTF{MISSING}
	{Object exists.}
	{Object MISSING is not present, as expected.}

\pxNewObject{CHK}
\pxSetValue{CHK}{nums}{1,2,3}
\pxIfObjectExistsTF{CHK}{CHK exists.}{CHK does not exist.}
\pxIfKeyExistsTF{CHK}{nums}{Key \texttt{nums} exists in CHK.}{Key missing.}
Length of CHK.nums: \pxListLen{CHK}{nums}.

\subsection{Boolean flags and scalar access}
\pxNewObject{FLAGS}
\pxSetValue{FLAGS}{on}{true}
\pxSetValue{FLAGS}{off}{{false}}
\pxIfTrueTF{FLAGS}{on}{ON}{OFF},\quad
\pxIfTrueTF{FLAGS}{off}{ON}{OFF}.

\pxNewObject{SCALAR}
\pxSetValue{SCALAR}{flag}{true}
\pxSetValue{SCALAR}{text}{Hello, world.}
Scalar SCALAR.flag: \pxGetScalar{SCALAR}{flag}.\\
Scalar SCALAR.text: \pxGetScalar{SCALAR}{text}.\\
Missing key: [\pxGetScalar{SCALAR}{none}].

\subsection{List clear and compact API}
\pxNewObject{LCLR}
\pxListSet{LCLR}{vals}{a,b,c}
Length before clear: \pxListLen{LCLR}{vals}.
\pxListClear{LCLR}{vals}
Length after clear: \pxListLen{LCLR}{vals}.
\pxListPutRight{LCLR}{vals}{x,y}
Values: \pxListAll{LCLR}{vals}.

\pxSETlst LCLR.vals={z,a}
After \verb|\pxSETlst|: \pxListAll{LCLR}{vals}.
\pxSET LCLR.vals=A
After \verb|\pxSET|: \pxGETtl(LCLR.vals).
\pxPUT LCLR.vals+=B
After \verb|\pxPUT|: \pxPRINTlst LCLR.vals().
With separator: \pxPRINTlst LCLR.vals(, ).

% ========== 3. pxUML — Class diagrams ==========
\section{pxUML: class diagrams}

\subsection{Classes and diagram}
\pxUMLClass{Person}
	{- name : String}
	[+ getName() : String]

\addAttribute{Person}{- age : Integer}
\addOperation{Person}{+ setName(newName : String) : void}

\pxUMLClass{BaseClass}{- id : Integer}
\pxUMLClass{DerivedClass}{- extra : String}
\addInheritance{DerivedClass}{BaseClass}

\pxUMLClass{Service}{- name : String}[+ run() : void]

\pxUMLDiagram{
	\drawclass[text width=14em]{Person}[0,0]
	\drawsimpleclass*[text width=10em]{BaseClass}[8,0]
	\drawsimpleclass[text width=10em]{DerivedClass}[8,-4]
	\drawinstance*[14]{ServiceInstance}[0,-6]{Service}{- name : String}{+ run() : void}
}[0.9\linewidth]

\PXshowclass{Person}
\PXshowclass{Service}

Inspection: Person attributes: [\pxUMLAttributes{Person}].\\
Service operations: [\pxUMLOperations{Service}].\\
Parents of DerivedClass: [\pxUMLParents{DerivedClass}].\\
Non-existent class: [\pxUMLAttributes{NoExisteix}].

\subsection{Enums and associations}
\pxUMLClass[enumclass]{Role}{DEVELOPER, MANAGER}
\pxUMLClass[enumclass]{Kind}{A, B, C}

\pxUMLClass{Entity}{- role : Role, - name : String}[+ doIt() : void]
\pxUMLClass{Container}{- id : Integer}

\pxUMLDiagram{
	\drawclass*[text width=12em]{Role}[0,0]
	\drawclass*[text width=12em]{Kind}[10,0]
	\drawclass[text width=14em]{Entity}[0,-4]
	\drawclass[text width=14em]{Container}[10,-4]
	\association{Entity}{0..*}{}{Container}{}{1}
}[0.85\linewidth]

% ========== 4. pxSRC — Images ==========
\section{pxSRC: safe images}

\pxIMG[.2\linewidth]{img/testimage}[primarycolor]{Single image caption}[fig:one]

\pxIMG[.2\linewidth]{img/testimage}{Caption (optional width; color uses default)}[fig:two]

\pxIMGpair{img/testimage}{img/testimage}[primarycolor]{Pair caption}[fig:pair]

\pxIMGList{testimage}{Step one, Step two, Step three, Step four}

% ========== 5. pxCORE — Debug (optional) ==========
\section{pxCORE: setup and debug}

\pxCOREsetup{debug=false}
Debug state: \pxCOREIfDebugTF{[DEBUG-ON]}{[DEBUG-OFF]}.

% ========== 6. pxGDX — Layout ==========
\section{pxGDX: layout helpers}

\pxGrow{Example title text}[\linewidth]

\end{document}
